\part{Supplementary Results}
% ==============================================================================
\section{Parametric equations}
% ------------------------------------------------------------------------------
\paragraph{Distribution families}
We considered 6 distribution families for $p(x)$ in our analysis
(Table~\ref{tab:fams}).
\begin{table}[h]
  \centering
  \caption{Distribution families considered,
    using common parameter notation ($\aa,\bb$)}
  \label{tab:fams}
  \newcommand{\dd}[4]{#1 & {$\displaystyle#4$} & $#2$ & $#3$\\[3ex]}
\begin{tabular}{llcc}
  \toprule
  Family & Definition $p(x)$ & $\aa$ log-prior\tn[*] & $\bb$ log-prior\tn[*] \\
  \midrule
  \dd{Exponential\tn[\dag]}        {   -   }{[-1,~3]}{\frac{1}{\aa}~e^{-\frac{x}{\aa} }}
  \dd{Gamma\tn[\dag]}              {[-2,~2]}{[-3,~3]}{\frac{\bb^\aa}{\Gamma(\aa)}~x^{\aa-1}e^{-\bb x}}
  \dd{Weibull\tn[\dag]}            {[-1,~2]}{[-1,~2]}{\frac{\aa}{\bb}~{\left(\frac{x}{\bb}\right)}^{\aa-1}e^{-{(x/\bb)}^\aa}}
  \dd{Log-Normal\tn[\dag]}         {[-9,~1]}{[-2,~3]}{\frac{1}{x~\bb\sqrt{2\pi}}~e^{-\frac12{\left(\frac{\log{x}-\aa}{\bb}\right)}^2}}
  \dd{Scaled Beta\tn[\ddag]}       {[-3,~3]}{[~0,~3]}{\frac{x^{\aa-1}{(1-x^\aa)}^{\bb-1}}{B(\aa,\bb)}}
  \dd{Scaled Kumaraswamy\tn[\ddag]}{[-2,~3]}{[~0,~3]}{\aa~\bb~x^{\aa-1}{(1-x^\aa)}^{\bb-1}}
  \\[-5ex]\bottomrule
\end{tabular}
\smallskip\floatfoot{
  \tnd[*]{Priors for $\aa,\bb$ in the applied model were defined as uniform on the log-scale;
    \eg $[-1,~3]$ corresponds to $[e^{-1},~e^{3}] \approx [0.37,~20]$.}\\
  \tnd[\dag] {Distributions spanning $[0,\infty)$ were truncated to $[0,50]$ years.}
  \tnd[\ddag]{Distributions spanning $[0,1]$ were scaled to span $[0,50]$ years.}}
% notation:
% exp:     aa = NA,          bb = 1/rate (β)
% gamma:   aa = shape   (α), bb = rate   (β)
% weibull: aa = shape   (k), bb = scale  (λ)
% lnorm:   aa = meanlog (?), bb = sdlog  (?)
% sbeta:   aa = shape1, (α), bb = shape2 (β)
% skumar:  aa = a       (a), bb = b      (b)

\end{table}
\clearpage % TEMP
% ------------------------------------------------------------------------------
\paragraph{Likelihoods functions}
Under the model for $p(z|s)$, we define the likelihoods
for empirical estimates of the mean $\sE_i[z|s]$ and proportions $\sP_i[\azb|s]$
from a given survey with sample size $N_i$ as, respectively:
the Gamma-equivalent of Normal sampling error for the mean (since $\sE > 0$), and
the Beta-approximation of the Binomial distribution (for differentiability):
\begin{alignat}{1}
  \sE_i[z|s;N_i] &\sim \txn{Gamma}(\aa = N_i~/~\CV_i^2[z|s],~~
                                   \bb = N_i~/~\CV_i^2[z|s]~/~\E[z|s])
  \label{eq:lE} \\
  \sP_i[\azb|s;N_i] &\sim \txn{Beta}(\aa = N_i~   \P_i[\azb|s],~~
                                     \bb = N_i~(1-\P_i[\azb|s]))
  \label{eq:lP}
\end{alignat}
\par
To validate these likelihood functions,
we simulated 10,000 stochastic samples of $z$
with sample sizes $N_i = \{30, 100, 300, 1000\}$
given known $p(z|s)$ under $\E[x] = 5, \CV[x] = \{0.5, 1, 1.5\}$,
and all 6 distribution families.
We then compared the distribution of stochastic estimates of
$\sE[z|s;N_i]$ and $\sP[\azb|s;N_i]$ from each sample
with the analytic likelihood functions: \eqrefa{eq:lE}{eq:lP}.
For proportions, we considered
$\P[\lzu01]$, $\P[\lzu15]$, and $\P[\lzu59]$, arbitrarily.
\par
Figures~\ref{fig:val.lf.mean}--\ref{fig:val.lf.p59} illustrate these comparisons,
where we observe good agreement between the
stochastic and analytic distributions across all conditions.
Note that stochastic distributions for $\sP[\azb|s;N_i]$
are discrete (more visible at low $N_i$) because
there are a finite number simulated respondents
whose durations can fall in the specified interval $[a,b]$.
\foreach \slug/\stat in {%
  mean/$\E[z|s]$,%
  p01/$\P[\lzu01|s]$,%
  p15/$\P[\lzu15|s]$,%
  p59/$\P[\lzu59|s]$}{
\begin{figure}
  \centering\includegraphics[width=\linewidth]{val.lf.\slug.pdf}
  \caption{Comparison of analytic likelihood functions
    and distribution of 10,000 stochastic estimates of \textbf{\stat},
    for different simulated survey sizes $N_i$ and
    latent distribution families for $p(x)$}
  \label{fig:val.lf.\slug}
\end{figure}}
\clearpage
% ==============================================================================
\section{Illustrative analytic results}
Figure~\ref{fig:toy.distr} repeats Figure~\ref{fig:toy.distr.1}
for 5 distribution families considered (Table~\ref{tab:fams}) --- \ie
showing distributions of total ($x$) and censored ($z$) durations
among the source, active, and sampled populations for each family,
all having $\E[x] = 5$ and $\CV[x] = \{0.5, 1, 1.5\}$.
We omit the Exponential family as it is
a special case of both Gamma and Weibull with $\CV[x] = 1$.
Figure~\ref{fig:toy.est} further illustrates
performance of different $\E[x]$ estimators considered in \cite{Fazito2012}
across a range of $\CV[x]$ and the same 5 distribution families.
\vspace{-2cm} % HACK
\begin{figure}
  \centering\includegraphics[width=\linewidth]{toy.distr}
  \caption{Illustrative distributions of total and censored durations
    among the source, active, and sampled populations for
    different $x$ distribution families (columns) and
    3 coefficients of variation $\CV[x]$ (rows)}
  \label{fig:toy.distr}
  \floatfoot{Note: true $\E[x] = 5$;
    probability of being sampled modelled per \eqref{eq:p(s|z)} with $\pp = 0.5, \tt = 1$.}
\end{figure}
\begin{figure}
  \centering\includegraphics[width=\linewidth]{toy.est}
  % TODO: investigate truncation issues @ CV->2
  \caption{Performance of $\E[x]$ estimators considered in \cite{Fazito2012}
    across a range of $\CV[x]$ and distribution families}
  \label{fig:toy.est}
  \floatfoot{Note: true $\E[x] = 5$;
    probability of being sampled modelled per \eqref{eq:p(s|z)} with $\pp = 0.5, \tt = 1$;
    means and medians for censored data ($z$) are doubled per \cite{Fazito2012}.}
\end{figure}
\clearpage
% ==============================================================================
\section{Application}
\def\postcap{Posterior distribution of
  mean \& CV total/censored distributions among source/sampled populations:
  $\E[x], \CV[x]~/~\E[z|s], \CV[z|s]$;
  model parameters: $\aa,\bb,\pp,\tt$;
  and log-likelihood: $\log(L)$
  across 6 distribution families considered}
% ------------------------------------------------------------------------------
\paragraph{Eswatini data}
Table~\ref{tab:data.demo} summarizes the
RDS-adjusted sex work duration data reported in \cite{Baral2014},
while Figure~\ref{fig:demo.pars} illustrates the posterior distributions
inferred by fitting our Bayesian model to those data.
\begin{table}[h]
  \centering
  \caption{RDS-adjusted proportions of durations selling sex
    among Swati female sex workers in \cite{Baral2014} (N = 313)}
  \label{tab:data.demo}
  \begin{tabular}{lrrcr}
\toprule
Duration\tn[\dag] & \multicolumn{3}{c}{Proportion (\%)} & $N_i$ \\
\cmidrule(rl){2-4}
(years)  & Crude & RDS & (\CI) & Eff.\tn[\ddag] \\
\midrule
0--3     & 28.8 & 38.3 & (27.5,~49.1) &  77 \\
3--6     & 31.3 & 32.1 & (23.6,~40.7) & 113 \\
6--11    & 26.2 & 20.2 & (13.2,~27.1) & 129 \\
11+      & 13.7 &  9.4 &  (4.4,~14.4) & 135 \\
\bottomrule
\end{tabular}
\smallskip\floatfoot{
  \tnd[\dag]{Reported strata in \cite{Baral2014} were:
    "under 2, 3--5, 6-10, 11 or more",
    which we corrected here to avoid implied gaps between strata},
  \tnd[\ddag]{Effective sample size per strata calculated
    to match RDS-adjusted \CI.}}

\end{table}
\begin{figure}[h]
  \centering\includegraphics[width=\textwidth]{demo.pars}
  \caption{\postcap\ 
    when fitting to the RDS-adjusted Eswatini data (Table~\ref{tab:data.demo})}
  \label{fig:demo.pars}
\end{figure}
\clearpage % TEMP
% ------------------------------------------------------------------------------
\paragraph{Global data}
\cite{Fazito2012}
% TODO: CSV file
